\documentclass[UTF8,11pt,a4paper]{ctexart}
\usepackage[a4paper,left=1.45cm,right=1.45cm,top=1.55cm,bottom=1.65cm]{geometry}
\usepackage{amsmath,amssymb,bm}
\usepackage{graphicx}
\usepackage{booktabs,tabularx,array,longtable}
\usepackage{enumitem}
\usepackage[dvipsnames,svgnames]{xcolor}
\usepackage{tikz}
\usetikzlibrary{arrows.meta,positioning,calc,fit,shapes.geometric}
\usepackage[most]{tcolorbox}
\usepackage{hyperref}
\usepackage{caption}
\usepackage{float}
\usepackage{fancyhdr}
\usepackage{titlesec}
\usepackage{microtype}
\IfFontExistsTF{Times New Roman}{\setmainfont{Times New Roman}}{\setmainfont{Liberation Serif}}
\IfFontExistsTF{Noto Sans CJK SC}{\setCJKmainfont{Noto Sans CJK SC}}{}
\IfFontExistsTF{Noto Sans CJK SC}{\setCJKsansfont{Noto Sans CJK SC}}{}
\IfFontExistsTF{DejaVu Sans Mono}{\setmonofont{DejaVu Sans Mono}}{\setmonofont{Latin Modern Mono}}
\definecolor{mainblue}{HTML}{0B5CAD}
\definecolor{softblue}{HTML}{EAF4FF}
\definecolor{darktext}{HTML}{1F2937}
\definecolor{softgray}{HTML}{F5F7FA}
\definecolor{accent}{HTML}{F59E0B}
\definecolor{goodgreen}{HTML}{2E7D32}
\definecolor{badred}{HTML}{B91C1C}
\definecolor{purple}{HTML}{6D5BD0}
\hypersetup{colorlinks=true,linkcolor=mainblue,urlcolor=DarkOrange,citecolor=mainblue}
\setlength{\parindent}{2em}
\setlength{\headheight}{14pt}
\setlength{\parskip}{0.22em}
\setlist[itemize]{leftmargin=2em,itemsep=0.23em,topsep=0.25em}
\setlist[enumerate]{leftmargin=2.2em,itemsep=0.23em,topsep=0.25em}
\renewcommand{\arraystretch}{1.2}
\captionsetup{font=small,labelfont=bf}
\titleformat{\section}{\Large\bfseries\sffamily\color{mainblue}}{\thesection}{0.7em}{}
\titleformat{\subsection}{\large\bfseries\sffamily\color{darktext}}{\thesubsection}{0.6em}{}
\titleformat{\subsubsection}{\normalsize\bfseries\sffamily\color{darktext}}{\thesubsubsection}{0.55em}{}
\pagestyle{fancy}
\fancyhf{}
\fancyhead[L]{\small\sffamily ChordEdit}
\fancyhead[R]{\small\sffamily \leftmark}
\fancyfoot[C]{\small\thepage}
\renewcommand{\headrulewidth}{0.4pt}
\newtcolorbox{infobox}[1]{enhanced,breakable,colback=softblue,colframe=mainblue!70!black,boxrule=0.8pt,arc=2mm,left=1.5mm,right=1.5mm,top=1mm,bottom=1mm,title=\bfseries #1}
\newtcolorbox{keybox}[1]{enhanced,breakable,colback=softgray,colframe=darktext!45,boxrule=0.6pt,arc=2mm,left=1.5mm,right=1.5mm,top=1mm,bottom=1mm,title=\bfseries #1}
\newtcolorbox{notebox}[1]{enhanced,breakable,colback=accent!8,colframe=accent!80!black,boxrule=0.7pt,arc=2mm,left=1.5mm,right=1.5mm,top=1mm,bottom=1mm,title=\bfseries #1}
\newcommand{\term}[1]{\textbf{\textcolor{mainblue}{#1}}}
\begin{document}

\section{论文基本信息}
\begin{center}
\begin{tabularx}{0.94\linewidth}{>{\bfseries}p{0.18\linewidth}X}
\toprule
题目 & ChordEdit\\
课件日期 & 2026-06-14\\
研究方向 & 图像编辑 / 最优传输\\
一句话概括 & 用最优传输和路径设计协调源图像、文本条件与目标分布,使编辑在保持原图结构的同时实现可靠语义变化。\\
\bottomrule
\end{tabularx}
\end{center}

\begin{infobox}{摘要要点}
用最优传输和路径设计协调源图像、文本条件与目标分布,使编辑在保持原图结构的同时实现可靠语义变化。 本说明按“研究问题—核心机制—基础公式—实验阅读—局限与优化”的顺序重新组织，不保留原始材料的封面、目录和页码对应。
\end{infobox}

\section{研究问题与动机}
这项工作的出发点不是简单地替换某个网络层，而是识别现有范式中最影响\term{质量、效率、可靠性或可部署性}的瓶颈。结合论文所讨论的问题，可以将研究动机概括为以下几点：
\begin{itemize}
\item 而ChordEdit “ ” 是顺推,顺着最自然、最省力的路径,把马慢慢 变成独角兽,整个过程平滑流畅,背景也不受任何影响
\item “ ” 弦路径:在坑洼的山地上修了一条高速公路,即使只用一步完 成编辑,不会出现剧烈的抖动。
\item ChordEdit :用最优传输理论,重新定义图像编辑
\item “ ” 前人的方法是硬掰,把马硬掰成独角兽,不符合逻辑
\item 最优传输理论早已存在,但少有人想到用在一步式图像编辑上
\item 当下论文的创新点用更多卡训了更大的模型达到更高的效果
\end{itemize}

\begin{notebox}{理解重点}
阅读这类论文时，应把“问题是否真实存在”“方法是否直接解决该问题”“实验是否排除了其他解释”分开判断。仅凭更大的模型、更多数据或更长训练得到的提升，不能自动视为结构创新带来的收益。
\end{notebox}

\section{核心思想与整体流程}
用最优传输和路径设计协调源图像、文本条件与目标分布,使编辑在保持原图结构的同时实现可靠语义变化。 从信息流角度看，其基本流程可以整理为下图。

\begin{figure}[H]
\centering
\begin{tikzpicture}[>=Latex,node distance=0.78cm]
\tikzstyle{procstep}=[draw,rounded corners,fill=softblue,align=center,minimum width=3.05cm,minimum height=0.95cm,text width=2.75cm]
\node[procstep] (a) {源图像与文本};
\node[procstep,right=of a] (b) {最优传输路径};
\node[procstep,right=of b] (c) {条件生成/编辑};
\node[procstep,right=of c] (d) {目标图像};
\draw[->,thick,mainblue] (a)--(b);
\draw[->,thick,mainblue] (b)--(c);
\draw[->,thick,mainblue] (c)--(d);
\end{tikzpicture}
\caption{核心信息流与处理阶段。}
\end{figure}

\subsection{方法组成与信息流}
\begin{itemize}
\item 做图像反演:先把真实图像映射回隐空间,找到对应的向量再编辑,但反演本身就会 丢失大量信息,编辑效果自然好不了
\item 额外训练编辑模型:专门训一个模型来修正变形,但成本极高,而且只能适配特定的 生成模型,失去了通用性
\item FlowEdit 需要33 步,运行时间长达7.22 秒
\item 最优传输理论早已存在,但少有人想到用在一步式图像编辑上
\item 当下论文的创新点用更多卡训了更大的模型达到更高的效果
\item 科研真正的魅力所在:用最简单的工具,解决最复杂的问题
\item 真正的科研从来都不是比谁的显卡多,而是比谁的脑子好
\end{itemize}

\begin{keybox}{结构解析}
\begin{itemize}
\item 粒子偏离目标分布三种解决方法 ； 增加推理步数：一步变成十步、二十步，但速度慢
\item ； FlowEdit 需要33 步，运行时间长达7.22 秒 ； 额外训练编辑模型：专门训一个模型来修正变形，但成本极高，而且只能适配特定的 生成模型，
\end{itemize}
\end{keybox}

上述内容以文字方式保留方法结构，避免把整张幻灯片作为独立页面嵌入正文。

\section{数学与机制理解}
本节给出理解该工作的基础数学结构。公式用于解释方法所属范式，并不替代原论文中的完整推导。
\begin{equation}
x_t=\alpha_t x_0+\sigma_t\varepsilon,\qquad \varepsilon\sim\mathcal N(0,I)
\end{equation}
扩散类方法把真实样本与噪声按时间混合，模型再学习逆向去噪、速度场或条件流。不同论文的主要差异通常落在表示空间、网络骨干和训练目标上。

\begin{keybox}{从公式回到方法}
应重点观察论文究竟改变了公式中的哪一部分：是输入表示、连接矩阵、条件变量、参数化方式、训练目标，还是求解与调度过程。真正的创新通常可以在这一层被准确定位。
\end{keybox}


\section{实验结果应如何阅读}
\begin{itemize}
\item ChordEdit :用最优传输理论,重新定义图像编辑
\item 实验结果:老泰坦跑出了吊打全场的性能
\item CVPR 2026 颁奖活动展示的算力消耗统计数据
\item 所有实验基于18 年Titan 显卡
\end{itemize}

\begin{notebox}{结果解读}
\begin{itemize}
\item 实验结果：老泰坦跑出了吊打全场的性能
\end{itemize}
\end{notebox}

下图用于直观呈现主要对比、消融、定性样例或效率结果。
\begin{figure}[H]
\centering
\includegraphics[width=0.94\linewidth,height=0.43\textheight,keepaspectratio]{result_evidence.jpg}
\caption{代表性实验对比、消融结果或定性样例。}
\end{figure}
实验部分不应只看单一最好数字。还需要核对模型规模、训练数据、训练步数、采样或解码预算、硬件条件以及是否使用额外蒸馏、检索或后处理。只有比较条件接近时，结果才能真正说明方法本身的优势。

\section{优点、局限与可优化空间}
\subsection{主要优点}
\begin{itemize}
\item \term{问题与方法对应明确}：论文围绕一个可识别的核心瓶颈展开，方法结构能够直接映射到该瓶颈。
\item \term{可复用的设计思想}：即使不完整复现模型，其中的分解、稀疏化、递归、条件控制或系统优化思想也可迁移到相邻任务。
\item \term{实验提供了多角度证据}：实验通常同时展示主结果、效率比较、案例或消融，使结论不完全依赖单一指标。
\end{itemize}

\subsection{局限与进一步优化}
\begin{itemize}
\item 需要进一步区分\term{结构收益}与更大算力、更多数据、不同训练技巧带来的收益，并在等参数、等计算预算下进行严格比较。
\item 对超参数、输入分布和模型规模的敏感性仍应系统分析；在一个数据集上的最优配置不一定能直接迁移。
\item 实际部署还要考虑显存峰值、并发、延迟抖动、失败案例和鲁棒性，而不仅是离线平均指标。
\item 可尝试将该方法与蒸馏、量化、缓存复用、动态路由或更强的表示学习结合，验证不同优化是否互补。
\end{itemize}

\section{结论}
总体而言，用最优传输和路径设计协调源图像、文本条件与目标分布，使编辑在保持原图结构的同时实现可靠语义变化。。 进一步需要注意：最优传输理论早已存在,但少有人想到用在一步式图像编辑上;当下论文的创新点用更多卡训了更大的模型达到更高的效果。

\end{document}
